% part: intuitionistic-logic
% chapter: propositions-as-types
% section: proofs-to-terms

\documentclass[../../../include/open-logic-section]{subfiles}

\begin{document}

\olfileid{int}{pty}{pt}

\olsection{تبدیل \usetoken{P}{derivation} به ترم‌های برهان}

فرایند تبدیل !!{proof}s استنتاج طبیعی به ترم‌های برهان، یعنی ترجمهٔ
!!{proof}s~\Log{N2i} به !!{proof}s~\Log{tN2i}، بسیار ساده است.
کافی است در سمت چپ هر !!{formula} واقع در راست یک دنباله در
!!{derivation}، ترم برهانی بنویسیم؛ در نتیجه دنباله‌هایی به شکل
$\Gamma \Sequent M : !A$ حاصل می‌شوند. سپس می‌گوییم $M$ در بافت
~$\Gamma$ \emph{گواه}~$!A$ است. قواعد \Log{tN2i} ترم برهانی را که
باید نوشته شود کاملاً تعیین می‌کنند.

\begin{prop}
هر !!{proof}~$\delta$ برای $\Gamma \Sequent !A$ در \Log{N2i} با
!!a{proof}~$\delta'$ برای $\Gamma \Sequent N : !A$ در \Log{tN2i}
متناظر است. ترم برهان~$N$ به‌طور یکتا با~$\delta$ تعیین می‌شود.
\end{prop}

\begin{proof}
بر ارتفاع~$\delta$ استقرا می‌کنیم. اگر $\delta$ اصل
$x:!A \Sequent !A$ باشد، $\delta'$ اصل~$x:!A \Sequent x:!A$ است.
اگر $\delta$ با قاعده‌ای استنتاجی پایان یابد، فرض استقرا برای هر مقدمه
ترم برهانی و برای دنبالهٔ شامل آن ترم، !!{proof}s در~\Log{tN2i} فراهم
می‌کند که با آن مقدمه متناظر است. قاعدهٔ متناظر \Log{tN2i} را اعمال
می‌کنیم؛ قاعدهٔ استنتاج \Log{tN2i} ترم برهان متناظر را تعیین می‌کند.
\end{proof}

\begin{ex}
  !!{proof} بسیار سادهٔ زیر را در \Log{N1i} برای
  $(!A \land !B) \lif (!B \land !A)$ در نظر بگیرید:
  \[
  \AxiomC{$\Discharge{!A\land !B}{x}$}
  \RightLabel{\Elim\land[2]}
  \UnaryInfC{$!B$}
  \AxiomC{$\Discharge{!A\land !B}{x}$}
  \RightLabel{\Elim\land[1]}
  \UnaryInfC{$!A$}
  \RightLabel{\Intro\land}
  \BinaryInfC{$!B \land !A$}
  \DischargeRule{\Intro\lif}{x}
  \UnaryInfC{$(!A \land !B) \lif (!B \land !A)$}
  \DisplayProof
  \]
  این برهان در \Log{N2i} چنین است:
  \[
  \Axiom$x : !A\land !B \fCenter !A\land !B$
  \RightLabel{\Elim\land[2]}
  \UnaryInf$x : !A\land !B \fCenter !B$
  \Axiom$x : !A\land !B \fCenter !A\land !B$
  \RightLabel{\Elim\land[1]}
  \UnaryInf$x : !A\land !B \fCenter !A$
  \RightLabel{\Intro\land}
  \BinaryInf$x : !A\land !B \fCenter !B \land !A$
  \RightLabel{\Intro\lif}
  \UnaryInf$\fCenter (!A \land !B) \lif (!B \land !A)$
  \DisplayProof
  \]
  ترم‌های برهان را از بالا به پایین اختصاص می‌دهیم. ترم‌های برهان
  اصول همان متغیر موجود در بافت، یعنی~$x$، هستند. ترم‌های برهان متناظر
  با $\Elim\land[i]$ سپس به صورت $\proj{2}{x}$ و~$\proj{1}{x}$ تعیین
  می‌شوند. با اعمال \Intro{\land} این دو ترم را ترکیب می‌کنیم:
  $\pair{\proj{2}{x}}{\proj{1}{x}}$. سرانجام، قاعدهٔ \Intro{\lif}
  یک انتزاع~$\lambd$ اعمال می‌کند، $x$ را مقید می‌سازد و ثبت می‌کند که
  $x$ فرض~$!A \land !B$ را برچسب زده بود:
  $\lambd[\typeof{x}{!A \land
  !B}][\pair{\proj{2}{x}}{\proj{1}{x}}]$.
  آنگاه !!{proof} در \Log{tN2i} چنین است:
  \[
  \Axiom$x : !A\land !B \fCenter x: !A\land !B$
  \RightLabel{\Elim\land[2]}
  \UnaryInf$x : !A\land !B \fCenter \proj{2}{x} : !B$
  \Axiom$x : !A\land !B \fCenter x: !A\land !B$
  \RightLabel{\Elim\land[1]}
  \UnaryInf$x : !A\land !B \fCenter \proj{1}{x} : !A$
  \RightLabel{\Intro\land}
  \BinaryInf$x : !A\land !B \fCenter \pair{\proj{2}{x}}{\proj{1}{x}} :
    !B \land !A$
  \RightLabel{\Intro\lif}
  \UnaryInf$\fCenter \lambd[\typeof{x}{!A \land
  !B}][\pair{\proj{2}{x}}{\proj{1}{x}}] : (!A \land !B) \lif (!B \land !A)$
  \DisplayProof
  \]
\end{ex}


\end{document}
